% Options for packages loaded elsewhere
\PassOptionsToPackage{unicode}{hyperref}
\PassOptionsToPackage{hyphens}{url}
%
\documentclass[
  10pt,
]{article}
\usepackage{amsmath,amssymb}
\usepackage{iftex}
\ifPDFTeX
  \usepackage[OT1]{fontenc}
  \usepackage[utf8]{inputenc}
\else % if luatex or xetex
  \usepackage{unicode-math} % this also loads fontspec
  \defaultfontfeatures{Scale=MatchLowercase}
  \defaultfontfeatures[\rmfamily]{Ligatures=TeX,Scale=1}
\fi
\ifPDFTeX\else
  % xetex/luatex font selection
\fi
% Use upquote if available, for straight quotes in verbatim environments
\IfFileExists{upquote.sty}{\usepackage{upquote}}{}
\IfFileExists{microtype.sty}{% use microtype if available
  \usepackage[]{microtype}
  \UseMicrotypeSet[protrusion]{basicmath} % disable protrusion for tt fonts
}{}
\makeatletter
\@ifundefined{KOMAClassName}{% if non-KOMA class
  \IfFileExists{parskip.sty}{%
    \usepackage{parskip}
  }{% else
    \setlength{\parindent}{0pt}
    \setlength{\parskip}{6pt plus 2pt minus 1pt}}
}{% if KOMA class
  \KOMAoptions{parskip=half}}
\makeatother
\usepackage{xcolor}
\usepackage[margin=.88in]{geometry}
\setlength{\emergencystretch}{3em} % prevent overfull lines
\providecommand{\tightlist}{%
  \setlength{\itemsep}{0pt}\setlength{\parskip}{0pt}}
\setcounter{secnumdepth}{-\maxdimen} % remove section numbering
\ifLuaTeX
  \usepackage{selnolig}  % disable illegal ligatures
\fi
\usepackage{bookmark}
\IfFileExists{xurl.sty}{\usepackage{xurl}}{} % add URL line breaks if available
\urlstyle{same}
\hypersetup{
  pdftitle={Fractional planks near the central equal-normal plane},
  pdfauthor={Leo Gao, AI-assisted research},
  hidelinks,
  pdfcreator={LaTeX via pandoc}}

\title{Fractional planks near the central equal-normal plane}
\author{Leo Gao, AI-assisted research}
\date{}

\begin{document}
\maketitle

\section{Sharp fractional planks on the central plane with equal
coefficients}\label{sharp-fractional-planks-on-the-central-plane-with-equal-coefficients}

% \textbf{Leo Gao, AI-assisted, private research note.}

\textbf{Theorem.} Let \(K\subseteq[0,1]^3\cap\{x_1+x_2+x_3=3/2\}\) be
convex and meet all six cube facets. Then every fractional coordinate
plank cover has cost at least \(2/3\). In fact \(K\) contains a
translated equilateral triangle of one of two fixed orientations whose
three coordinate projections all have width \(2/3\), supporting a law
uniform on each full projection interval. No permutation symmetry,
central symmetry, or packed triangle is assumed. Equivalently, a convex
subset of \([-1,1]^3\cap\{x_1+x_2+x_3=0\}\) meeting all facets supports
a Lebesgue-dominated matching of mass \(4/3\). This does not assert that
\(2/3\) is attained within the central subclass, nor settle other
planes.

\subsection{The geometric assertion}\label{the-geometric-assertion}

The ambient hexagon has vertices the permutations of \((0,1/2,1)\).
Translate by its center and write its consecutive vertices
\(v_0,\ldots,v_5\). In a planar affine coordinate system they are \[
(1,0), (1/2,1), (-1/2,1), (-1,0), (-1/2,-1), (1/2,-1).
\] An original cube coordinate face cuts out exactly one hexagon edge.
Choose a contact \(p_i=(1-s_i)v_i+s_i v_{i+1}\) on each edge, where
\(0\le s_i\le1\); indices are mod six. The hull of these six points lies
in \(K\).

The six vertices have the identity \(v_{i-1}+v_{i+1}=v_i\). Thus the
chord between \(p_{i-1}\) and \(p_i\) contains \((1-h_i)v_i\), where \[
h_i=\frac{a_i b_i}{a_i+b_i},\quad a_i=1-s_{i-1},\quad b_i=s_i.
\] When \(a_i+b_i=0\), set \(h_i=0\). Indeed the convex weights
\(b_i/(a_i+b_i)\) and \(a_i/(a_i+b_i)\) give exactly the indicated
point. Since every \(h_i\) is at most 1/2, the chord points on opposite
rays show that the origin is in the hull. In particular
\((2/3)v_i\in K\) whenever \(h_i\le1/3\), using the origin. Call a
vertex \emph{bad} when \(h_i>1/3\), meaning this particular radial
construction is insufficient (not implying that the point is absent from
\(K\)). Since \(a_i,b_i\le1\), a bad vertex requires \(a_i,b_i>1/2\).
Equivalently it has \(s_{i-1}<1/2<s_i\) along with the sharper
inequality on \(h_i\). Hence no adjacent vertices are bad.

If all bad vertices have the same parity, or there are none, the three
vertices of the other parity already give the required triangle.
Otherwise there is exactly one bad vertex of each parity, and they are
opposite: the distance between distinct nonadjacent vertices of opposite
parity in a six-cycle is three, so after choosing one of either parity,
the other is forced. Relabel the pair as \(v_0,v_3\), choosing \(v_0\)
with smaller deficit \(\delta=h_0-1/3\le h_3-1/3=:\eta\). Both lie in
\((0,1/6]\), because every \(h_i\le1/2\). We prove \[
(2/3)v_i-\delta v_0\in K\qquad(i=0,2,4).
\] For \(i=0\) this is the chord point \((1-h_0)v_0\). For \(i=2\) it is
the point with coordinates \((-1/3-\delta,2/3)\); for \(i=4\) reflect
about the horizontal axis and use the same argument below.

Write \(s=s_2\in[0,1/2)\) and \(r=s_3\in(1/2,1]\), according to badness
of \(v_3\). Here \[
p_2=(-1/2-s/2,1-s),\qquad p_3=(-1+r/2,-r).
\] If \(s\le1/3\), the segment from \(p_2\) to \(p_3\) reaches height
\(2/3\) with horizontal coordinate at most \(-1/2\le-1/3-\delta\). If
\(1/3<s<1/2\), instead use the segment from \(p_2\) to
\(p_1=(1/2-s_1,1)\). At height \(2/3\) its horizontal coordinate is at
most \[
 1/2-\frac{1+s/2}{3s}=1/3-\frac1{3s}=-1/3-\frac{1-2s}{3s}.
\] The last margin is at least \(\eta\): since \(r\le1\), \[
\eta=\frac{(1-s)r}{1-s+r}-\frac13
   \le\frac{1-s}{2-s}-\frac13
   =\frac{1-2s}{3(2-s)}
   \le\frac{1-2s}{3s}.
\] Hence the height \(2/3\) section has a point far enough left in both
cases.

It has a point far enough right as well. If \(s_0\ge2/3\), the segment
between \(p_5\) and \(p_0\) reaches height \(2/3\) and has horizontal
coordinate at least \(1/2\). If \(s_0<2/3\), the segment between
\(p_0=(1-s_0/2,s_0)\) and \(p_1\) reaches that height at horizontal
coordinate at least \[
 -1/2+\frac{1}{3(1-s_0)}(3/2-s_0/2)
 =\frac{s_0}{3(1-s_0)}\ge0.
\] Convexity now gives the target point. Reflection in the horizontal
axis sends \(v_i\) to \(v_{-i}\), interchanges \(a_3,b_3\), hence
preserves \(\eta\), and fixes \(v_0\) and the shift direction. Thus the
same calculation yields the \(i=4\) target. This finishes the geometric
assertion, including contact endpoints and equality cases of the
auxiliary inequalities.

\subsection{Uniform projection
certificate}\label{uniform-projection-certificate}

Undo the affine display coordinates. Each parity of three hexagon
vertices consists of cyclic permutations of \((0,1/2,-1/2)\) or its
negative after translation back to the true central plane coordinates.
After multiplication by \(2/3\) the centered coordinates in a parity
triple are permutations of \((0,1/3,-1/3)\); all three projection widths
are \(2/3\). The same holds after any common translation, including the
constructed shift. Put probability \(1/3\) uniformly on each triangle
edge. In any coordinate, before the translation, the three projected
intervals are \([-1/3,0]\), \([0,1/3]\), and \([-1/3,1/3]\). The
resulting density is \(3/2\) on the whole interval. Rescaling the law to
mass \(2/3\) gives Lebesgue domination; integrating the fractional cover
proves the theorem.

The earlier symmetry-based result needed cyclic invariance of the entire
body. This proof uses only the rotation symmetry of the ambient plane
section; the six chosen contact parameters are arbitrary. The general
planar AHKM problem is still open in this work.

\subsection{Independent exact stress
tests}\label{independent-exact-stress-tests}

The accompanying \texttt{verify-central-hexagon.py} verifies each
constructed triangle by a rational convex-hull calculation, rather than
reusing the proof's inequalities. It passes all 15,625 contact
configurations on the quarter-edge grid (including 1,200 with opposite
bad corners), plus 2,000 exactly rational seeded configurations (194
opposite-bad). The analytic proof above establishes the full continuum;
the tests are cross-checks, not its basis.

\section{A strict gap on the central equal-normal plane, stable under
perturbation}\label{a-strict-gap-on-the-central-equal-normal-plane-stable-under-perturbation}

% \textbf{Leo Gao, AI-assisted, private supplementary research note.}

\textbf{Theorem.} There exist \(\varepsilon>0\) and \(\eta>0\) with the
following property. Let \(K\subseteq[0,1]^3\) be any packed convex set
in an affine plane \(a_1x_1+a_2x_2+a_3x_3=t\), normalized by
\(\sum a_i=3\). If \[
\max_i |a_i-1|+|t-3/2|<\varepsilon,
\] then \(\tau^*(K)\ge 2/3+\eta\). No symmetry or Hausdorff-nearness of
\(K\) to the \emph{complete} section is assumed. The constants here are
nonexplicit. In particular this extends the bound to an open
neighborhood of the central equal-normal plane in both normal and
offset, not just to one affine plane.

The proof refines the geometry in \href{central-hexagon.md}{the
central-hexagon note}, then uses a compactness argument with a density
regularization that controls perturbations of the projection directions.

\subsection{1. Pointwise strictness for central contact
hulls}\label{pointwise-strictness-for-central-contact-hulls}

Use the notation \(v_i,p_i,s_i,h_i\) from the central-hexagon proof. Let
\(H\) be the convex hull of the six chosen contacts. The radial points
\((1-h_i)v_i\in H\) have positive radii, in all six directions, so the
origin belongs to the interior of \(H\) (for instance, \(H\) contains
the hull of the points \(v_i/2\)). Call a corner \emph{strictly bad} if
\(h_i>1/3\) and \emph{tight} if \(h_i=1/3\). A radially good corner with
\(h_i<1/3\) has \((2/3)v_i\) in the interior of \(H\): it belongs to the
open segment from the interior origin to \((1-h_i)v_i\).

\textbf{Case A: there are strictly bad corners of both parities.} There
is exactly one of each, opposite, relabeled \(0,3\) with
\(0<\delta=h_0-1/3\le\eta_0=h_3-1/3\le1/6\). The shifted even triangle
from the central proof has vertices \((2/3)v_i-\delta v_0\) for
\(i=0,2,4\). If \(\delta<1/6\), the side vertices \(i=2,4\) in that
proof are interior: the horizontal section has a point strictly on
either side of the proposed coordinate. Indeed, the left bound \(-1/2\)
in the first case is strictly below \(-1/3-\delta\), and in the second
case \[
\frac{1-2s}{3s}>\frac{1-2s}{3(2-s)}\ge\eta_0\ge\delta
\quad (1/3<s<1/2).
\] The right bounds are strictly positive relative to the target. The
horizontal heights \(\pm2/3\) lie between the hull's extremes \(-1,1\),
attained on the top and bottom edges. Thus a point strictly inside such
a horizontal section is interior in the plane: a supporting line through
it would have to contain that section and be horizontal. The third
target \((1-h_0)v_0\) becomes interior if we translate a little farther
along \(-v_0\), staying on the open radial segment. Both side vertices
remain interior for a sufficiently small additional translation. Then
enlarge the triangle a little to obtain common coordinate width strictly
above \(2/3\).

If \(\delta=1/6\), then \(h_0=h_3=1/2\). Equality in \(ab/(a+b)\le1/2\)
for \(a,b\in[0,1]\) requires \(a=b=1\). Hence the hull contains the four
vertices \(v_1,v_2,v_4,v_5\). In true cube coordinates centered at
\((1/2,1/2,1/2)\), after relabeling or global reflection,
\(v_1=(1/2,0,-1/2)\), \(v_2=(0,1/2,-1/2)\), \(v_4=-v_1\), \(v_5=-v_2\).
Put equal probability on the uniform laws of \([v_1,-v_2]\) and
\([-v_1,v_2]\). Each of the first two coordinates is uniform on one
half-interval in each law and gets the two opposite halves in the
mixture. Each law already has third coordinate uniform on
\([-1/2,1/2]\). Thus the fractional capacity is 1 in this case.

\textbf{Case B: strictly bad corners occur in at most one parity.} A
tight corner must have \(a_i,b_i\ge1/2\). If any corner is strictly bad,
each adjacent corner has either \(a\) or \(b\) strictly below \(1/2\),
so is not tight. Choose the other parity, with no strictly bad corners
and at most one tight one (the possible opposite corner). If there are
no strictly bad corners, choose a parity with at most two tight ones;
otherwise all six would be tight, forcing every \(s_i=1/2\) by the
inequalities \(s_i\ge1/2\) and \(s_{i-1}\le1/2\) at every corner. At
\(s_i=1/2\) all \(h_i=1/4\), a contradiction.

For each tight vertex \(i\), the target \((2/3)v_i\) lies strictly
between \(p_{i-1}\) and \(p_i\); the points are distinct because both
\(a_i,b_i\ge1/2\). In local display coordinates with \(v_i=(1,0)\), the
line through the two endpoints has equation \[
 x=2/3+\kappa y,\qquad\kappa=1/(3b_i)-1/2,\quad |\kappa|\le1/6.
\] This follows by substitution of \(p_i=(1-b_i/2,b_i)\). The origin is
strictly to the left. In the regular Euclidean metric, with height
coordinate \(\sqrt3\,y/2\), the line's outward unit normal makes angle
at most \(\arctan(1/(3\sqrt3))<30^\circ\) with the radial direction
\(v_i\). Two vertices of one parity are separated by \(120^\circ\), so
the corresponding two outward normals, if needed, are not antiparallel.
One can choose a common strict inward displacement \(w\): with one tight
vertex use minus its normal, and with two use minus the sum of their
unit normals. For each tight vertex the displacement moves the point
into the interior of the triangle
\(\operatorname{conv}\{0,p_{i-1},p_i\}\) for sufficiently small positive
multiples of \(w\), since it starts in the relative interior of the edge
and crosses strictly inward. The other vertices, already interior to
\(H\), remain there. Enlarge slightly after this translation.

Thus for every central contact configuration, a Lebesgue-dominated
measure of mass strictly greater than \(2/3\) is available, either from
the edge law on an enlarged motif or from the explicit parallelogram
law.

\subsection{2. Stability under varying the affine
plane}\label{stability-under-varying-the-affine-plane}

Suppose no uniform constants existed. For each \(k\) choose packed
\(K_k\) in a plane whose normalized coefficients and offset differ from
the central ones by less than \(1/k\) in the displayed quantity, and
whose covering number is less than \(2/3+1/k\). Choose six facet
contacts from each \(K_k\); their compact hull is \(H_k\subseteq K_k\).
Pass to a subsequence on which the six contacts converge. Their limits
are the corresponding central hexagon contacts. Let \(H\) be their hull.
By Part 1, it is full-dimensional in its plane and supports a
Lebesgue-dominated measure \(\mu\) of mass \(m>2/3\).

Identify the central plane with coordinates \((x_1,x_2)\). The perturbed
plane maps from this display space by \[
 \Phi_k(u,v)=(u,v,b_k-c_k u-d_k v),\qquad b_k\to3/2,\quad c_k,d_k\to1.
\] The pullbacks of \(H_k\) converge in planar Hausdorff distance to the
display hull of \(H\), by convergence of the six generators. First
contract \(\mu\) toward an interior point of \(H\) by a factor \(s<1\)
chosen close enough to one that \(sm>2/3\), and multiply its mass by
\(s\). All marginal density bounds remain one. Then convolve in the
plane with a nonnegative smooth probability kernel supported in a
sufficiently small Euclidean disk. This keeps all three marginal density
bounds one: projecting a convolution is a one-dimensional convolution by
a probability, which preserves the bound. The resulting measure has mass
\(m'=sm>2/3\) and smooth compactly supported planar density \(g\), with
support a positive distance inside \(H\). The existence of that margin
follows from convexity: if \(c+RB\subset H\) is an interior ball and
\(s<1\), the contracted hull plus \((1-s)RB\) is still contained in
\(H\).

A fixed compact subset of the interior of a convex hull remains in all
sufficiently Hausdorff-close convex hulls: use support functions and the
positive ball margin. Thus for large \(k\) the support of \(g\) maps
under \(\Phi_k\) inside \(H_k\). The first two projected densities are
unchanged. The third is \[
 h_k(z)=\frac1{c_k}\int_{\mathbb R}g\!\left(\frac{b_k-z-d_k v}{c_k},v\right)\,dv,
\] with \(c_k>0\) eventually. These densities converge uniformly to the
central density, since \(g\) is uniformly continuous with compact
support, \(v\) can be restricted to a fixed compact interval, and all
\(h_k\) vanish outside a fixed bounded interval. Hence all three
marginal densities have sup norm at most \(1+o(1)\). Rescale the mass by
\(1/(1+o(1))\) using a nonnegative error bound to produce a
Lebesgue-dominated matching on \(H_k\) of mass tending to \(m'>2/3\).
This contradicts the claimed covering numbers and proves the uniform
theorem.

This compactness transfer covers genuinely unequal normal coefficients
and offsets near the middle of the cube. It does not establish a
quantitative size for the neighborhood or give the sharp bound for all
planar packed bodies.

\end{document}
